%% ============================================================
%%  compare_equipo_06: comparación de familias y de validaciones
%%  Plantilla Beamer Uniandes (beamerthemeUniandes.sty, sin cambios).
%%  Compilar con LuaLaTeX:  latexmk -lualatex compare_equipo_06.tex
%%  Las cifras y tablas salen de generated/ (scripts/36_build_presentations.py);
%%  las figuras, de ../views/ (scripts 32 y 33). Notas del orador: notas_compare.md.
%% ============================================================
\documentclass[aspectratio=169,11pt]{beamer}
%% Preamble: packages and theme. Keep design changes in beamerthemeUniandes.sty
\usepackage[utf8]{inputenc}
\usepackage[T1]{fontenc}
\usepackage[spanish,es-nodecimaldot,es-noshorthands]{babel}
\usepackage{amsmath,amssymb}
\usepackage{graphicx}
\usepackage{booktabs}
\usepackage[round]{natbib}
\usetheme{Uniandes}

\graphicspath{{../views/}}
% Generado por scripts/36_build_presentations.py a partir de views/. No editar.
\newcommand{\KagBest}{145,7}
\newcommand{\KagPrivate}{148,6}
\newcommand{\RankPublic}{1}
\newcommand{\RankPrivate}{1}
\newcommand{\SharePublic}{20\,\%}
\newcommand{\SharePrivate}{80\,\%}
\newcommand{\KagSecond}{150,8}
\newcommand{\GapSecond}{2,2}
\newcommand{\PrivateRise}{2,9}
\newcommand{\PrivateRisePct}{2,0\,\%}
\newcommand{\GainNaive}{61\,\%}
\newcommand{\KagBoost}{152,1}
\newcommand{\KagNet}{157,1}
\newcommand{\KagLinear}{180,5}
\newcommand{\KagScaled}{171,7}
\newcommand{\KagQsixtyfive}{161,7}
\newcommand{\KagText}{153,1}
\newcommand{\GainText}{7,9}
\newcommand{\GainBetoKaggle}{2,6}
\newcommand{\GainFirst}{216,1}
\newcommand{\GainLast}{15,3}
\newcommand{\GainTotal}{231,4}
\newcommand{\KagStepA}{377,1}
\newcommand{\KagStepB}{180,1}
\newcommand{\KagStepC}{160,9}
\newcommand{\KagStepD}{152,1}
\newcommand{\KagStepE}{148,2}
\newcommand{\KagStepF}{145,7}
\newcommand{\GainStepB}{196,9}
\newcommand{\GainStepC}{19,2}
\newcommand{\GainStepD}{8,9}
\newcommand{\GainStepE}{3,8}
\newcommand{\GainStepF}{2,6}
\newcommand{\MaeBeto}{148,5}
\newcommand{\MaeRidge}{169,1}
\newcommand{\BetoCorr}{0,83}
\newcommand{\BetoTopDecile}{341}
\newcommand{\RidgeTopDecile}{537}
\newcommand{\BetoMidDecile}{133}
\newcommand{\RidgeMidDecile}{87}
\newcommand{\BetoTopBias}{−314}
\newcommand{\RidgeTopBias}{−534}
\newcommand{\SpatialBoostingBefore}{127,0}
\newcommand{\SpatialBoostingAfter}{124,4}
\newcommand{\SpatialRedBefore}{137,6}
\newcommand{\SpatialRedAfter}{133,7}
\newcommand{\SpatialMezclaBefore}{127,6}
\newcommand{\SpatialMezclaAfter}{124,8}
\newcommand{\SpatialBest}{124,8}
\newcommand{\RandomBefore}{95,5}
\newcommand{\TwinsBefore}{112,2}
\newcommand{\RecipeGainMin}{12,3}
\newcommand{\RecipeGainMax}{44,8}
\newcommand{\LinSpatial}{164,1}
\newcommand{\LinNoGeoSpatial}{156,3}
\newcommand{\CorrForest}{0,98}
\newcommand{\CorrNet}{0,96}
\newcommand{\BlendNetGain}{−3,8}
\newcommand{\BlendForestGain}{0,0}
\newcommand{\BlendForestSpatial}{+2,8}
\newcommand{\NNegative}{16}
\newcommand{\MaxNet}{2.002}
\newcommand{\MaxForest}{1.590}
\newcommand{\MaxCart}{1.580}
\newcommand{\OsmSpatial}{+7,9 a +8,6}
\newcommand{\LocalSpatial}{+8,9 a +60,3}
\newcommand{\BalanceSpatial}{−7,7 a +5,5}
\newcommand{\BalanceWeighted}{+0,7 a +4,1}
\newcommand{\StackWeighted}{+15,1 a +23,2}
\newcommand{\SpRandom}{0,44}
\newcommand{\SpTwins}{0,78}
\newcommand{\SpSpatial}{0,74}
\newcommand{\SpWeighted}{0,83}
\newcommand{\GapSpatial}{26,5}
\newcommand{\AsymWinner}{190}
\newcommand{\AsymWinnerQ}{0,20}
\newcommand{\AsymWinnerCost}{174}
\newcommand{\BestCostKthree}{289}
\newcommand{\BestRankKthree}{6}
\newcommand{\BestRankKone}{4}
\newcommand{\NRanked}{13}
\newcommand{\OverClosed}{67\,\%}
\newcommand{\OverCapital}{18\,\%}
\newcommand{\OverMean}{123}
\newcommand{\LowClosed}{33\,\%}
\newcommand{\LowCapital}{12\,\%}
\newcommand{\LowMean}{67}
\newcommand{\SameDiscount}{15\,\%}
\newcommand{\SameClosed}{34\,\%}
\newcommand{\SameCapital}{14\,\%}
\newcommand{\MinCapital}{10\,\%}
\newcommand{\AptClosed}{32\,\%}
\newcommand{\AptCapital}{10,9\,\%}
\newcommand{\HouseCapital}{14,9\,\%}
\newcommand{\ShareOverCheap}{92\,\%}
\newcommand{\ShareOverRich}{18\,\%}
\newcommand{\BiasRich}{−205}
\newcommand{\ShareOverTotal}{67\,\%}
\newcommand{\BiasTotal}{39,4}
\newcommand{\MeanOver}{123,3}
\newcommand{\MeanUnder}{127,9}
\newcommand{\OverHighPred}{204}
\newcommand{\MaeApt}{108,3}
\newcommand{\MaeHouse}{175,7}
\newcommand{\MaeAreaOk}{117,0}
\newcommand{\MaeAreaMissing}{141,7}
\newcommand{\ShareAreaMissing}{26\,\%}
\newcommand{\BiasNortheast}{−53,3}


\title[Comparación de modelos y validaciones]{Por qué gana la mezcla de boosting y red neuronal}
\subtitle{Comparación de familias y de esquemas de validación}
\author[Equipo 06]{[NOMBRES]}
\institute{Big Data y Machine Learning (MECA 4107)\\ Universidad de los Andes}
\date[Nov.\ 2026]{Bogotá, noviembre de 2026}

\newcommand{\twolines}[2]{\begin{tabular}[t]{@{}c@{}}#1\\#2\end{tabular}}
% Figura a casi todo el ancho y, debajo, las viñetas que dicen qué mirar en ella.
\newcommand{\figslide}[2][0.88]{\centering\includegraphics[width=#1\textwidth]{#2}\par\vspace{2pt}}
\newenvironment{lookat}{\begin{itemize}\small\setlength{\itemsep}{1pt}}{\end{itemize}}

\begin{document}

{\setbeamertemplate{footline}{}
 \begin{frame}[plain,noframenumbering]
   \titlepage
 \end{frame}}

% 1 ---------------------------------------------------------------
\begin{frame}{La mezcla de boosting y red quedó primera en los dos tableros}
  \vfill
  \begin{columns}[T]
    \column{0.32\textwidth}\centering \keyfig{\KagBest\ M}{\twolines{tablero público}{(\SharePublic\ del test)}}
    \column{0.32\textwidth}\centering \keyfig{\KagPrivate\ M}{\twolines{tablero privado, 1.\textsuperscript{er} puesto}{(\GapSecond\ M sobre el segundo)}}
    \column{0.32\textwidth}\centering \keyfig{\SpatialBest\ M}{\twolines{validación espacial: zonas}{que no vio al entrenar}}
  \end{columns}
  \vfill
  \begin{itemize}\small
    \item MAE: error promedio en pesos entre el precio predicho y el observado.
    \item Boosting: muchos árboles pequeños; cada uno corrige el error del anterior.
    \item Red neuronal: modelo que puede predecir por encima de los precios que vio.
    \item Mezcla: promedio ponderado de los dos, 70\,\% boosting y 30\,\% red.
  \end{itemize}
  \vfill
\end{frame}

% 2 ---------------------------------------------------------------
\begin{frame}{El boosting gana la espacial y el tablero; la red, el ponderado}
  \centering\scriptsize
  % Generado por scripts/36_build_presentations.py. No editar.
\begin{tabular}{lrlrrrr}
\toprule
Familia & Envío & Búsqueda & Repetidos & Espacial & Ponderado & Público \\
\midrule
Regresión lineal & 108 & No aplica (MCO) & 144,1 & 164,1 & 166,0 & 180,5 \\
Elastic net & 104 & Optuna TPE (15) & 144,5 & 158,5 & 165,8 & sin envío \\
Elastic net + boosting & 49 & -- & -- & 137,4 & 162,3 & 184,1 \\
Elastic net + boosting q0,65 & 54 & -- & -- & 142,7 & 156,4 & 173,3 \\
CART & 105 & Optuna TPE (25) & 141,3 & 159,8 & 168,7 & 177,5 \\
Random forest & 114 & Optuna TPE (25) & 123,1 & 143,8 & 159,6 & 158,9 \\
Boosting & 127 & Probadas, sin mejora & 110,2 & 127,0 & 155,5 & -- \\
Boosting & 83 & Probadas, sin mejora & 111,0 & 131,5 & 154,8 & 152,1 \\
Red neuronal & 148 & Probadas, sin mejora & 124,8 & 137,6 & 148,2 & 157,1 \\
Mezcla & 152 & Pesos fijos 70/30 & 112,2 & 127,6 & 151,7 & 148,2 \\
\textbf{Mezcla con BETO} & \textbf{203} & \textbf{Pesos fijos 70/30} & \textbf{--} & \textbf{124,8} & \textbf{--} & \textbf{145,7} \\
\bottomrule
\end{tabular}

  \par\vspace{2pt}{\scriptsize\color{andesgray} MAE en millones de pesos; menos es mejor. La elastic net no se subió a Kaggle.}
  \vfill
  \begin{itemize}\small
    \item Gemelos: validación aleatoria (anuncios al azar de zonas ya vistas), con las copias juntas.
    \item Ponderado: holdout al lado de Chapinero, con los pesos del test.
    \item La red gana el ponderado, pero en el tablero queda detrás del boosting.
    \item Sin ponderado para la 203: BETO se afinó con precios de ese holdout.
  \end{itemize}
  \vfill
\end{frame}

% 3 ---------------------------------------------------------------
\begin{frame}{Los árboles no predicen por encima del tope del train; la red sí}
  \figslide{family_test_range}
  \begin{lookat}
    \item Línea punteada: el precio más alto del train, 1.650 M.
    \item CART (\MaxCart\ M) y bosque (\MaxForest\ M) quedan bajo el tope: un árbol promedia lo que vio.
    \item La red 148 llega a \MaxNet\ M: por eso aporta en los inmuebles caros.
  \end{lookat}
\end{frame}

% 4 ---------------------------------------------------------------
\begin{frame}{La red mejora la mezcla; el bosque comete los mismos errores}
  \figslide{family_blend_with_127}
  \begin{lookat}
    \item Cada punto: mezcla 70/30 del boosting 127 con otra familia; abajo es mejor.
    \item Bosque: correlación de errores de \CorrForest; el ponderado cambia \BlendForestGain\ M.
    \item Red: correlación de \CorrNet\ y buena sola; el ponderado cambia \BlendNetGain\ M.
  \end{lookat}
\end{frame}

% 5 ---------------------------------------------------------------
\begin{frame}{La receta común mejora a todas las familias}
  \figslide{common_recipe}
  \begin{lookat}
    \item Receta: por tipo, precio en pesos, cuantil 0,70, estrato, terreno y ridge del texto.
    \item Cuantil 0,70: apuntar a un precio que el 70\,\% de los parecidos no supera.
    \item Todas ganan entre \RecipeGainMin\ y \RecipeGainMax\ M en el ponderado; en la espacial, ninguna pasa al 83.
  \end{lookat}
\end{frame}

% 6 ---------------------------------------------------------------
\begin{frame}{BETO mejora a los dos miembros de la mezcla}
  \vfill
  \centering\footnotesize
  % Generado por scripts/36_build_presentations.py. No editar.
\begin{tabular}{lcrrr}
\toprule
Modelo & Filas & Sin BETO & Con BETO & Cambio \\
\midrule
Boosting & 127 → 198 & 127,0 & 124,4 & −2,5 \\
Red neuronal & 148 → 196 & 137,6 & 133,7 & −3,9 \\
\textbf{Mezcla 70/30} & \textbf{152 → 203} & \textbf{127,6} & \textbf{124,8} & \textbf{−2,8} \\
\bottomrule
\end{tabular}

  \par\vspace{2pt}{\scriptsize\color{andesgray} Validación espacial, MAE en millones de pesos.}
  \vfill
  \begin{itemize}\small
    \item Ridge sobre TF-IDF: regresión que le pone un peso a cada palabra del anuncio.
    \item BETO: modelo de lenguaje en español, afinado aquí para estimar el precio desde el texto.
    \item Con BETO junto al ridge, la mezcla baja \GainBetoKaggle\ M en el tablero público.
    \item Solo lo probamos en el boosting y la red, no en las demás familias.
  \end{itemize}
  \vfill
\end{frame}

% 7 ---------------------------------------------------------------
\begin{frame}{Todos los esquemas son optimistas frente al tablero}
  \figslide[0.8]{validation_vs_kaggle_wide}
  \begin{lookat}
    \item Validación aleatoria: probar en anuncios al azar, de las mismas zonas que ya vio.
    \item Puntos sobre la diagonal: el esquema promete menos error que el tablero público.
    \item Spearman (1 = mismo orden que el tablero): aleatoria \SpRandom, espacial \SpSpatial.
  \end{lookat}
\end{frame}

% 8 ---------------------------------------------------------------
\begin{frame}{Los anuncios repetidos inflan la validación aleatoria}
  \figslide{validation_twins}
  \begin{lookat}
    \item Círculo vacío: aleatoria simple; lleno: gemelos en la misma partición; raya: tablero.
    \item Con los gemelos juntos, el MAE de la mezcla 152 sube de \RandomBefore\ a \TwinsBefore\ M.
    \item Aun así, todos quedan lejos del tablero: Chapinero es otra zona.
  \end{lookat}
\end{frame}

% 9 ---------------------------------------------------------------
\begin{frame}{La validación espacial premia bajar el nivel; el tablero no}
  \figslide{validation_level}
  \begin{lookat}
    \item Izquierda: con menos sesgo (cuantil más bajo), el MAE espacial baja.
    \item El ponderado casi no cambia entre los cuantiles 0,60 y 0,70.
    \item Derecha: el 0,65 gana en la espacial, pero no en el tablero (\KagQsixtyfive\ frente a \KagStepC\ M).
  \end{lookat}
\end{frame}

% 10 --------------------------------------------------------------
\begin{frame}{Ninguna de las \NNegative\ variantes mejoró los dos criterios a la vez}
  \figslide{negative_results}
  \begin{lookat}
    \item Cada punto: una variante frente a su referencia; a la derecha o arriba es peor.
    \item La esquina de abajo a la izquierda, que mejora ambas, está vacía.
    \item OpenStreetMap empeora la espacial (\OsmSpatial\ M): en Chapinero casi todo queda cerca.
  \end{lookat}
\end{frame}

% 11 --------------------------------------------------------------
\begin{frame}{Cada familia gana en algo y paga en otra cosa}
  \begin{itemize}
    \item Boosting: el mejor en zonas que no vio, pero no pasa el tope del train.
    \item Red neuronal: pasa el tope, pero sola es peor en el tablero (\KagNet\ M).
    \item Regresión lineal: fácil de leer, pero viaja mal a Chapinero (\KagLinear\ M).
    \item BETO: lee mejor el texto, pero exige GPU y su ponderado no es limpio.
  \end{itemize}
  \vfill
  \takeaway{El mejor modelo para el tablero no es el mejor para comprar: si sobrepagar cuesta el triple, la 203 cae al puesto \BestRankKthree.}
\end{frame}

\andesclosing{¡Gracias!}{Equipo 06 · MECA 4107}

\end{document}
